\documentclass[../DoAn.tex]{subfiles}
\begin{document}

Chương này tổng hợp kết quả benchmark cuối cùng của hệ thống mô phỏng. File dữ liệu gốc có 97 lần chạy; báo cáo chỉ lấy lần chạy mới nhất của mỗi \texttt{test\_name}, tương ứng với 64 kịch bản duy nhất. Cách lọc này tránh việc một kịch bản bị tính nhiều lần sau các lượt hiệu chỉnh.

\section{Dữ liệu và chỉ số đánh giá}

Mỗi lần chạy ghi lại cấu hình bản đồ, năng lượng tối đa, tỷ lệ bao phủ, số bước, năng lượng tiêu thụ, số lần quay về, số lần sạc, trạng thái kết thúc và đường dẫn ảnh chụp cuối. Các chỉ số chính được tóm tắt trong Bảng~\ref{tab:evaluation_metrics}.

\begin{table}[!htbp]
\centering
\caption{Các chỉ số sử dụng trong đánh giá thực nghiệm}
\label{tab:evaluation_metrics}
\small
\setlength{\tabcolsep}{4pt}
\renewcommand{\arraystretch}{1.08}
\begin{tabular}{|L{0.25\textwidth}|L{0.65\textwidth}|}
\hline
\textbf{Chỉ số} & \textbf{Ý nghĩa} \\
\hline
Coverage & Tỷ lệ ô hợp lệ đã được bao phủ khi kết thúc mô phỏng. \\
\hline
Steps & Tổng số bước di chuyển rời rạc của robot. \\
\hline
Energy & Tổng năng lượng chuẩn hóa đã tiêu thụ. \\
\hline
Ret. / Rec. & Số lần quay về căn cứ và số lần sạc lại. \\
\hline
Base & Robot có kết thúc tại căn cứ hay không. \\
\hline
Outcome & Trạng thái kết thúc: \texttt{SUCCESS}, \texttt{PARTIAL} hoặc \texttt{FAILED}. \\
\hline
Dyn. & Số ô bị ảnh hưởng bởi vật cản động trong kịch bản. \\
\hline
\end{tabular}
\end{table}

Trong các bảng sau, tên kịch bản được rút gọn thành mã như \texttt{A01}, \texttt{G10} hoặc \texttt{H17\_D14} để tránh tràn bảng. Tên đầy đủ vẫn được giữ trong file benchmark.

\section{Thiết kế bộ kiểm thử}

Bộ kiểm thử được chia thành tám nhóm. Bảng~\ref{tab:scenario_groups} cho biết biến số chính của từng nhóm. Các nhóm A, B, C, E, F và G đánh giá môi trường tĩnh. Nhóm D đưa vật cản động vào quá trình mô phỏng. Nhóm H đóng vai trò đối chứng cho mô hình toàn hướng lý tưởng, trong đó chi phí quay được bỏ qua ở mức lập kế hoạch.

\begin{table}[!htbp]
\centering
\caption{Cấu trúc bộ kịch bản kiểm thử cuối cùng}
\label{tab:scenario_groups}
\small
\setlength{\tabcolsep}{4pt}
\renewcommand{\arraystretch}{1.08}
\begin{tabular}{|C{0.10\textwidth}|L{0.24\textwidth}|L{0.56\textwidth}|}
\hline
\textbf{Nhóm} & \textbf{Kịch bản} & \textbf{Yếu tố được kiểm tra} \\
\hline
A & A01--A10 & Vị trí căn cứ trên cùng một bản đồ. \\
\hline
B & B01--B09 & Dung lượng năng lượng tối đa. \\
\hline
C & C01--C10 & Mật độ vật cản và tính liên thông của bản đồ. \\
\hline
D & D13--D14 & Vật cản động ảnh hưởng tới đường đi hoặc đường quay về. \\
\hline
E & E01--E10 & Chi phí địa hình có trọng số. \\
\hline
F & F01--F05 & Kích thước bản đồ. \\
\hline
G & G01--G11 & Cấu trúc phân bố vật cản khi mật độ gần như giữ nguyên. \\
\hline
H & H01--H06, H17\_D14 & Mô hình toàn hướng lý tưởng dùng làm đối chứng. \\
\hline
\end{tabular}
\end{table}

\section{Kết quả theo vị trí căn cứ}

Nhóm A giữ nguyên bản đồ chính kích thước 35$\times$30 và thay đổi vị trí căn cứ. Kết quả trong Bảng~\ref{tab:base_position_results} cho thấy tất cả kịch bản đều đạt 100\% coverage và kết thúc thành công.

\begin{table}[!htbp]
\centering
\caption{Kết quả nhóm A khi thay đổi vị trí căn cứ}
\label{tab:base_position_results}
\small
\setlength{\tabcolsep}{3pt}
\renewcommand{\arraystretch}{1.08}
\resizebox{\textwidth}{!}{%
\begin{tabular}{|C{0.08\textwidth}|R{0.10\textwidth}|R{0.09\textwidth}|R{0.11\textwidth}|R{0.07\textwidth}|R{0.07\textwidth}|R{0.09\textwidth}|C{0.10\textwidth}|}
\hline
\textbf{KC} & \textbf{Coverage} & \textbf{Steps} & \textbf{Energy} & \textbf{Ret.} & \textbf{Rec.} & \textbf{Remain} & \textbf{Outcome} \\
\hline
\texttt{A01} & 100.00\% & 1370 & 1574.5 & 9 & 8 & 32.5 & \texttt{SUCCESS} \\
\hline
\texttt{A02} & 100.00\% & 1816 & 2055.5 & 12 & 11 & 88 & \texttt{SUCCESS} \\
\hline
\texttt{A03} & 100.00\% & 1704 & 1925 & 11 & 10 & 42 & \texttt{SUCCESS} \\
\hline
\texttt{A04} & 100.00\% & 1638 & 1848.5 & 11 & 10 & 91.5 & \texttt{SUCCESS} \\
\hline
\texttt{A05} & 100.00\% & 1516 & 1722 & 10 & 9 & 64 & \texttt{SUCCESS} \\
\hline
\texttt{A06} & 100.00\% & 1524 & 1733 & 10 & 9 & 45 & \texttt{SUCCESS} \\
\hline
\texttt{A07} & 100.00\% & 1494 & 1698.5 & 10 & 9 & 75 & \texttt{SUCCESS} \\
\hline
\texttt{A08} & 100.00\% & 1534 & 1755 & 10 & 9 & 33.5 & \texttt{SUCCESS} \\
\hline
\texttt{A09} & 100.00\% & 1386 & 1589.5 & 9 & 8 & 0.5 & \texttt{SUCCESS} \\
\hline
\texttt{A10} & 100.00\% & 1394 & 1606.5 & 9 & 8 & 12 & \texttt{SUCCESS} \\
\hline
\end{tabular}%
}
\end{table}

Vị trí căn cứ vẫn tạo ra khác biệt đáng kể về chi phí. A01 dùng 1574.5 đơn vị năng lượng, trong khi A02 cần 2055.5 đơn vị. Sự khác biệt này phù hợp với vai trò của căn cứ trong mô hình: căn cứ không chỉ là điểm xuất phát mà còn là điểm robot phải quay về khi kiểm tra ràng buộc năng lượng.

\section{Ảnh hưởng của năng lượng tối đa}

Nhóm B giảm dần năng lượng tối đa trên cùng một bản đồ nhỏ. Bảng~\ref{tab:energy_capacity_results} cho thấy hệ thống vẫn hoàn thành toàn bộ các kịch bản.

\begin{table}[!htbp]
\centering
\caption{Kết quả nhóm B khi thay đổi năng lượng tối đa}
\label{tab:energy_capacity_results}
\small
\setlength{\tabcolsep}{3pt}
\renewcommand{\arraystretch}{1.08}
\resizebox{\textwidth}{!}{%
\begin{tabular}{|C{0.07\textwidth}|R{0.08\textwidth}|R{0.10\textwidth}|R{0.08\textwidth}|R{0.10\textwidth}|R{0.06\textwidth}|R{0.06\textwidth}|R{0.08\textwidth}|C{0.10\textwidth}|}
\hline
\textbf{KC} & \textbf{Emax} & \textbf{Coverage} & \textbf{Steps} & \textbf{Energy} & \textbf{Ret.} & \textbf{Rec.} & \textbf{Remain} & \textbf{Outcome} \\
\hline
\texttt{B01} & 260 & 100.00\% & 176 & 198 & 1 & 0 & 62 & \texttt{SUCCESS} \\
\hline
\texttt{B02} & 220 & 100.00\% & 176 & 198 & 1 & 0 & 22 & \texttt{SUCCESS} \\
\hline
\texttt{B03} & 180 & 100.00\% & 194 & 217 & 2 & 1 & 142 & \texttt{SUCCESS} \\
\hline
\texttt{B04} & 150 & 100.00\% & 216 & 241 & 2 & 1 & 59 & \texttt{SUCCESS} \\
\hline
\texttt{B05} & 120 & 100.00\% & 180 & 204 & 2 & 1 & 34 & \texttt{SUCCESS} \\
\hline
\texttt{B06} & 100 & 100.00\% & 206 & 231 & 3 & 2 & 62 & \texttt{SUCCESS} \\
\hline
\texttt{B07} & 70 & 100.00\% & 216 & 243 & 4 & 3 & 32 & \texttt{SUCCESS} \\
\hline
\texttt{B08} & 60 & 100.00\% & 260 & 291 & 5 & 4 & 4 & \texttt{SUCCESS} \\
\hline
\texttt{B09} & 50 & 100.00\% & 336 & 374 & 8 & 7 & 1 & \texttt{SUCCESS} \\
\hline
\end{tabular}%
}
\end{table}

Khi năng lượng tối đa giảm, số chu kỳ quay về và sạc lại tăng. B01 và B02 chỉ cần một lần quay về cuối nhiệm vụ. B09, với mức năng lượng tối đa 50, cần tám lần quay về và bảy lần sạc. Điều này cho thấy chính sách năng lượng không chỉ kiểm tra khả năng đi tới mục tiêu trước mắt mà còn duy trì điều kiện quay về căn cứ.

\section{Mật độ vật cản và tính liên thông}

Nhóm C thay đổi mật độ vật cản và tách riêng hai trường hợp: bản đồ còn liên thông và bản đồ bị chia cắt. Kết quả được trình bày trong Bảng~\ref{tab:obstacle_density_results}.

\begin{table}[!htbp]
\centering
\caption{Kết quả nhóm C theo mật độ vật cản}
\label{tab:obstacle_density_results}
\small
\setlength{\tabcolsep}{3pt}
\renewcommand{\arraystretch}{1.08}
\resizebox{\textwidth}{!}{%
\begin{tabular}{|C{0.07\textwidth}|R{0.08\textwidth}|R{0.10\textwidth}|R{0.08\textwidth}|R{0.10\textwidth}|R{0.06\textwidth}|R{0.06\textwidth}|C{0.06\textwidth}|C{0.10\textwidth}|}
\hline
\textbf{KC} & \textbf{Obs.} & \textbf{Coverage} & \textbf{Steps} & \textbf{Energy} & \textbf{Ret.} & \textbf{Rec.} & \textbf{Base} & \textbf{Outcome} \\
\hline
\texttt{C01} & 4.95\% & 100.00\% & 1430 & 1594 & 7 & 6 & Có & \texttt{SUCCESS} \\
\hline
\texttt{C02} & 10.00\% & 100.00\% & 1424 & 1615.5 & 7 & 6 & Có & \texttt{SUCCESS} \\
\hline
\texttt{C03} & 15.05\% & 100.00\% & 1522 & 1776 & 8 & 7 & Có & \texttt{SUCCESS} \\
\hline
\texttt{C04} & 20.00\% & 100.00\% & 1516 & 1816 & 8 & 7 & Có & \texttt{SUCCESS} \\
\hline
\texttt{C05} & 24.95\% & 100.00\% & 1540 & 1887 & 8 & 7 & Có & \texttt{SUCCESS} \\
\hline
\texttt{C06} & 30.00\% & 100.00\% & 1656 & 2045 & 9 & 8 & Có & \texttt{SUCCESS} \\
\hline
\texttt{C07} & 35.05\% & 100.00\% & 2882 & 3577 & 15 & 14 & Có & \texttt{SUCCESS} \\
\hline
\texttt{C08} & 40.00\% & 100.00\% & 2080 & 2618.5 & 11 & 10 & Có & \texttt{SUCCESS} \\
\hline
\texttt{C09} & 35.05\% & 93.40\% & 1630 & 2067.5 & 9 & 9 & Có & \texttt{PARTIAL} \\
\hline
\texttt{C10} & 40.00\% & 78.41\% & 1198 & 1518 & 7 & 7 & Có & \texttt{PARTIAL} \\
\hline
\end{tabular}%
}
\end{table}

Các kịch bản C01--C08 đều đạt \texttt{SUCCESS}, kể cả khi mật độ vật cản tăng lên 40\% trong trường hợp bản đồ vẫn liên thông. C09 và C10 kết thúc \texttt{PARTIAL} vì một phần ô hợp lệ nằm ngoài vùng robot có thể tiếp cận. Đây là kết quả mong muốn: robot quay về được căn cứ, nhưng báo cáo không xem nhiệm vụ là hoàn thành đầy đủ khi còn vùng không thể bao phủ.

\section{Cấu trúc phân bố vật cản}

Nhóm G bổ sung các bản đồ có mật độ vật cản gần như giữ nguyên nhưng hình dạng phân bố khác nhau. Bảng~\ref{tab:obstacle_structure_results} dùng để kiểm tra liệu hệ thống chỉ nhạy với mật độ vật cản hay còn chịu ảnh hưởng bởi cấu trúc hình học của bản đồ.

\begin{table}[!htbp]
\centering
\caption{Kết quả nhóm G theo cấu trúc phân bố vật cản}
\label{tab:obstacle_structure_results}
\small
\setlength{\tabcolsep}{3pt}
\renewcommand{\arraystretch}{1.08}
\resizebox{\textwidth}{!}{%
\begin{tabular}{|C{0.07\textwidth}|R{0.08\textwidth}|R{0.10\textwidth}|R{0.08\textwidth}|R{0.10\textwidth}|R{0.06\textwidth}|R{0.06\textwidth}|C{0.10\textwidth}|}
\hline
\textbf{KC} & \textbf{Free} & \textbf{Coverage} & \textbf{Steps} & \textbf{Energy} & \textbf{Ret.} & \textbf{Rec.} & \textbf{Outcome} \\
\hline
\texttt{G01} & 840 & 100.00\% & 1300 & 1569.5 & 4 & 3 & \texttt{SUCCESS} \\
\hline
\texttt{G02} & 840 & 100.00\% & 1254 & 1449 & 4 & 3 & \texttt{SUCCESS} \\
\hline
\texttt{G03} & 840 & 100.00\% & 1216 & 1387 & 4 & 3 & \texttt{SUCCESS} \\
\hline
\texttt{G04} & 840 & 100.00\% & 1172 & 1361 & 3 & 2 & \texttt{SUCCESS} \\
\hline
\texttt{G05} & 840 & 100.00\% & 1250 & 1466 & 4 & 3 & \texttt{SUCCESS} \\
\hline
\texttt{G06} & 840 & 100.00\% & 1264 & 1445.5 & 4 & 3 & \texttt{SUCCESS} \\
\hline
\texttt{G07} & 840 & 100.00\% & 1380 & 1606 & 4 & 3 & \texttt{SUCCESS} \\
\hline
\texttt{G08} & 840 & 100.00\% & 1344 & 1554 & 4 & 3 & \texttt{SUCCESS} \\
\hline
\texttt{G09} & 840 & 100.00\% & 1284 & 1459 & 4 & 3 & \texttt{SUCCESS} \\
\hline
\texttt{G10} & 840 & 100.00\% & 1098 & 1213.5 & 3 & 2 & \texttt{SUCCESS} \\
\hline
\texttt{G11} & 858 & 100.00\% & 1022 & 1102.5 & 3 & 2 & \texttt{SUCCESS} \\
\hline
\end{tabular}%
}
\end{table}

Tất cả kịch bản nhóm G đều đạt \texttt{SUCCESS}. Tuy vậy, tổng năng lượng thay đổi rõ. G11 dùng 1102.5 đơn vị năng lượng, trong khi G07 dùng 1606 đơn vị. Do nhiều kịch bản trong nhóm G có cùng mật độ 20\%, kết quả này cho thấy mật độ vật cản không đủ để mô tả độ khó của bản đồ. Hình dạng vùng trống, cổ hẹp, cụm vật cản và cách mở của vùng chữ U đều có thể làm thay đổi đường bao phủ.

\section{Vật cản động}

Nhóm D kiểm tra tình huống đường đi hoặc đường quay về bị ảnh hưởng bởi vật cản động. Bảng~\ref{tab:dynamic_obstacle_results} trình bày kết quả hai kịch bản đại diện.

\begin{table}[!htbp]
\centering
\caption{Kết quả nhóm D với vật cản động}
\label{tab:dynamic_obstacle_results}
\small
\setlength{\tabcolsep}{3pt}
\renewcommand{\arraystretch}{1.08}
\resizebox{\textwidth}{!}{%
\begin{tabular}{|C{0.08\textwidth}|R{0.10\textwidth}|R{0.08\textwidth}|R{0.10\textwidth}|R{0.07\textwidth}|C{0.07\textwidth}|C{0.10\textwidth}|}
\hline
\textbf{KC} & \textbf{Coverage} & \textbf{Steps} & \textbf{Energy} & \textbf{Dyn.} & \textbf{Base} & \textbf{Outcome} \\
\hline
\texttt{D13} & 100.00\% & 136 & 150 & 2 & Không & \texttt{PARTIAL} \\
\hline
\texttt{D14} & 98.59\% & 128 & 146 & 4 & Không & \texttt{FAILED} \\
\hline
\end{tabular}%
}
\end{table}

D13 đạt 100\% coverage nhưng không kết thúc tại căn cứ, nên được ghi là \texttt{PARTIAL}. D14 đạt 98.59\% coverage và kết thúc \texttt{FAILED}. Hai trường hợp này cho thấy coverage cao chưa đủ để kết luận nhiệm vụ thành công. Trạng thái cuối phải xét thêm vị trí robot, năng lượng còn lại và khả năng quay về căn cứ.

\section{Địa hình có trọng số}

Nhóm E thay đổi chi phí địa hình trên bản đồ. Kết quả trong Bảng~\ref{tab:weighted_terrain_results} cho thấy tất cả kịch bản đều hoàn thành, nhưng năng lượng tiêu thụ khác nhau rất mạnh.

\begin{table}[!htbp]
\centering
\caption{Kết quả nhóm E với bản đồ có trọng số địa hình}
\label{tab:weighted_terrain_results}
\small
\setlength{\tabcolsep}{3pt}
\renewcommand{\arraystretch}{1.08}
\resizebox{\textwidth}{!}{%
\begin{tabular}{|C{0.07\textwidth}|R{0.10\textwidth}|R{0.08\textwidth}|R{0.10\textwidth}|R{0.06\textwidth}|R{0.06\textwidth}|R{0.08\textwidth}|C{0.10\textwidth}|}
\hline
\textbf{KC} & \textbf{Coverage} & \textbf{Steps} & \textbf{Energy} & \textbf{Ret.} & \textbf{Rec.} & \textbf{Remain} & \textbf{Outcome} \\
\hline
\texttt{E01} & 100.00\% & 1118 & 1159.5 & 3 & 2 & 100 & \texttt{SUCCESS} \\
\hline
\texttt{E02} & 100.00\% & 1308 & 2383.5 & 6 & 5 & 130 & \texttt{SUCCESS} \\
\hline
\texttt{E03} & 100.00\% & 1506 & 3031 & 8 & 7 & 287 & \texttt{SUCCESS} \\
\hline
\texttt{E04} & 100.00\% & 1332 & 1950.5 & 5 & 4 & 146 & \texttt{SUCCESS} \\
\hline
\texttt{E05} & 100.00\% & 1538 & 4333 & 11 & 10 & 256 & \texttt{SUCCESS} \\
\hline
\texttt{E06} & 100.00\% & 1928 & 3025 & 8 & 7 & 306 & \texttt{SUCCESS} \\
\hline
\texttt{E07} & 100.00\% & 1584 & 4993.5 & 12 & 11 & 5 & \texttt{SUCCESS} \\
\hline
\texttt{E08} & 100.00\% & 1484 & 4090.5 & 10 & 9 & 86.5 & \texttt{SUCCESS} \\
\hline
\texttt{E09} & 100.00\% & 1254 & 1521.5 & 4 & 3 & 156 & \texttt{SUCCESS} \\
\hline
\texttt{E10} & 100.00\% & 1852 & 6128.5 & 15 & 14 & 114 & \texttt{SUCCESS} \\
\hline
\end{tabular}%
}
\end{table}

E01 là bản đồ phẳng, tiêu thụ 1159.5 đơn vị năng lượng. E10 tiêu thụ 6128.5 đơn vị và cần 15 lần quay về. Như vậy, số ô hợp lệ và kích thước bản đồ không đủ để dự đoán chi phí nhiệm vụ; trọng số địa hình là thành phần cần thiết của mô hình.

\section{Khả năng mở rộng theo kích thước bản đồ}

Nhóm F tăng kích thước bản đồ từ 20$\times$20 đến 40$\times$35. Kết quả được trình bày trong Bảng~\ref{tab:scale_results}.

\begin{table}[!htbp]
\centering
\caption{Kết quả nhóm F khi tăng kích thước bản đồ}
\label{tab:scale_results}
\small
\setlength{\tabcolsep}{3pt}
\renewcommand{\arraystretch}{1.08}
\resizebox{\textwidth}{!}{%
\begin{tabular}{|C{0.07\textwidth}|C{0.10\textwidth}|R{0.08\textwidth}|R{0.10\textwidth}|R{0.08\textwidth}|R{0.10\textwidth}|R{0.06\textwidth}|R{0.06\textwidth}|C{0.10\textwidth}|}
\hline
\textbf{KC} & \textbf{Size} & \textbf{Free} & \textbf{Coverage} & \textbf{Steps} & \textbf{Energy} & \textbf{Ret.} & \textbf{Rec.} & \textbf{Outcome} \\
\hline
\texttt{F01} & 20\$\textbackslash{}times\$20 & 320 & 100.00\% & 558 & 679 & 4 & 3 & \texttt{SUCCESS} \\
\hline
\texttt{F02} & 25\$\textbackslash{}times\$25 & 500 & 100.00\% & 836 & 1019.5 & 4 & 3 & \texttt{SUCCESS} \\
\hline
\texttt{F03} & 30\$\textbackslash{}times\$30 & 720 & 100.00\% & 1152 & 1387 & 4 & 3 & \texttt{SUCCESS} \\
\hline
\texttt{F04} & 30\$\textbackslash{}times\$35 & 840 & 100.00\% & 1406 & 1694 & 5 & 4 & \texttt{SUCCESS} \\
\hline
\texttt{F05} & 35\$\textbackslash{}times\$40 & 1120 & 100.00\% & 1922 & 2317.5 & 6 & 5 & \texttt{SUCCESS} \\
\hline
\end{tabular}%
}
\end{table}

Số ô hợp lệ tăng từ 320 ở F01 lên 1120 ở F05. Tương ứng, số bước tăng từ 558 lên 1922 và năng lượng tăng từ 679.0 lên 2317.5. Hệ thống vẫn giữ được trạng thái \texttt{SUCCESS} ở toàn bộ nhóm F, cho thấy quy trình lập kế hoạch, quay về và sạc lại vẫn ổn định trong phạm vi bản đồ được kiểm thử.

\section{So sánh với mô hình toàn hướng lý tưởng}

Nhóm H bỏ qua chi phí quay để tạo mô hình đối chứng. Bảng~\ref{tab:omni_results} liệt kê các kịch bản H cuối cùng.

\begin{table}[!htbp]
\centering
\caption{Kết quả nhóm H với mô hình toàn hướng lý tưởng}
\label{tab:omni_results}
\small
\setlength{\tabcolsep}{3pt}
\renewcommand{\arraystretch}{1.08}
\resizebox{\textwidth}{!}{%
\begin{tabular}{|C{0.10\textwidth}|R{0.10\textwidth}|R{0.08\textwidth}|R{0.10\textwidth}|R{0.06\textwidth}|R{0.06\textwidth}|R{0.06\textwidth}|C{0.06\textwidth}|C{0.10\textwidth}|}
\hline
\textbf{KC} & \textbf{Coverage} & \textbf{Steps} & \textbf{Energy} & \textbf{Ret.} & \textbf{Rec.} & \textbf{Dyn.} & \textbf{Base} & \textbf{Outcome} \\
\hline
\texttt{H01} & 100.00\% & 1204 & 1204 & 7 & 6 & 0 & Có & \texttt{SUCCESS} \\
\hline
\texttt{H02} & 100.00\% & 1594 & 1594 & 9 & 8 & 0 & Có & \texttt{SUCCESS} \\
\hline
\texttt{H03} & 100.00\% & 1346 & 1346 & 7 & 6 & 0 & Có & \texttt{SUCCESS} \\
\hline
\texttt{H04} & 100.00\% & 208 & 208 & 2 & 1 & 0 & Có & \texttt{SUCCESS} \\
\hline
\texttt{H05} & 100.00\% & 180 & 180 & 2 & 1 & 0 & Có & \texttt{SUCCESS} \\
\hline
\texttt{H06} & 100.00\% & 192 & 192 & 2 & 1 & 0 & Có & \texttt{SUCCESS} \\
\hline
\texttt{H17\_D14} & 100.00\% & 148 & 148 & 1 & 0 & 4 & Có & \texttt{SUCCESS} \\
\hline
\end{tabular}%
}
\end{table}

Ở cặp A01--H01, số bước giảm từ 1370 xuống 1204, tương đương giảm 12.12\%. Năng lượng giảm từ 1574.5 xuống 1204, tức giảm 23.53\%. Kết quả tương tự xuất hiện ở cặp A02--H02 và A10--H03. Điều này xác nhận rằng chi phí quay ảnh hưởng trực tiếp tới hành vi lập kế hoạch, không chỉ là phần cộng thêm sau khi đã chọn đường.

Cặp D14--H17\_D14 cần đọc cẩn thận hơn. D14 thất bại trước khi hoàn thành nhiệm vụ, còn H17\_D14 đạt \texttt{SUCCESS}. Vì trạng thái kết thúc khác nhau, không nên so sánh năng lượng của hai dòng như hai nghiệm tương đương. Điểm đáng chú ý là mô hình toàn hướng tạo điều kiện thuận lợi hơn trong một bản đồ có vật cản động và hành lang hẹp.

\section{Nhận xét tổng hợp}

Bộ benchmark cuối cùng cho thấy hệ thống đạt mục tiêu mô phỏng đã đặt ra. Trên bản đồ tĩnh liên thông, robot bao phủ toàn bộ vùng hợp lệ và quay về căn cứ. Khi năng lượng giảm, hệ thống tăng số lần quay về thay vì tiếp tục di chuyển thiếu an toàn. Khi bản đồ bị chia cắt, kết quả chuyển sang \texttt{PARTIAL} thay vì báo thành công giả. Khi có vật cản động, trạng thái cuối phản ánh đúng rủi ro không quay về được căn cứ. Các nhóm E và G bổ sung thêm bằng chứng rằng chi phí địa hình và cấu trúc vật cản đều ảnh hưởng tới năng lượng, số bước và số chu kỳ làm việc.

Các kết quả này không nhằm khẳng định thuật toán là tối ưu toàn cục cho mọi bài toán Coverage Path Planning. Ý nghĩa chính của thực nghiệm là kiểm chứng sự phối hợp giữa bốn thành phần: Dijkstra có hướng, chi phí địa hình, ràng buộc năng lượng quay về và cơ chế xử lý môi trường thay đổi. Đây là cơ sở để Chương 6 tổng kết kết quả và nêu hướng phát triển tiếp theo.

\end{document}
